% ---------------------------------------------------------------------------
% Capítulo 1 — Descripción de la empresa
% ---------------------------------------------------------------------------
% La norma 3.13 lo exige de forma explícita: «el cuerpo del trabajo debe incluir
% [...] Descripción de la empresa (para el caso de la Pasantía Larga/Intermedia)».
% Mínimo dos páginas por capítulo (norma 2.6). Redactado en tercera persona.
%
% Este capítulo recoge únicamente información verificada. Los datos proceden de
% la carta de aceptación CCT-001-2026 y de la información pública de la empresa.
%
% ESTE CAPÍTULO YA CUMPLE LA NORMA. Comprobado el 25/08/2026: la 3.13 solo
% exige que el cuerpo incluya una «Descripción de la empresa», sin enumerar
% contenidos, y la 2.6 pide un mínimo de dos páginas — el capítulo ocupa tres.
%
% Lo que sigue es OPCIONAL: enriquecería el capítulo si el tutor industrial lo
% facilita, pero no es requisito y no bloquea nada. La lista procede de las
% DIRECTRICES TSU-2024, que son del Decanato de Estudios Tecnológicos y NO
% aplican a esta pasantía.
%
% OPCIONAL   - Reseña histórica y año de constitución de la compañía.
% OPCIONAL   - Tamaño de la plantilla y su distribución por áreas.
% OPCIONAL   - Misión y visión, si están formalmente declaradas.
% OPCIONAL   - Organigrama o relación de áreas funcionales.
% OPCIONAL   - Equipo concreto de adscripción durante la pasantía.
% Queda anotado como punto 2, prioridad baja, en docs/cambios_pendientes.md.

\chapter{Descripción de la empresa}

La pasantía se desarrolla en Domotai C.A., una empresa venezolana de desarrollo
de software cuya actividad se concentra en la aplicación de la inteligencia
artificial a productos y procesos. Este capítulo describe su actividad, sus
líneas de servicio y su forma de trabajo, y sitúa el proyecto de la pasantía
dentro de ese contexto.

\section{Identificación y ubicación}

La empresa opera bajo la razón social Domotai C.A. y tiene su sede en el Edificio
Caracas Campus, calle Altagracia, urbanización La Trinidad, municipio Baruta del
estado Miranda, en la ciudad de Caracas.

Su ámbito de operación no se limita al mercado local: los proyectos referidos en
su portafolio corresponden a clientes de sectores diversos, atendidos bajo un
modelo de trabajo que no exige presencia física del cliente.

\section{Actividad y servicios}

Domotai C.A. se define como una empresa de desarrollo de software que combina el
criterio de sus profesionales con herramientas de inteligencia artificial para la
construcción de productos. Su oferta abarca las líneas recogidas en la tabla
\ref{tab:servicios}.

\begin{table}[H]
  \centering
  \caption{Líneas de servicio de Domotai C.A.}
  \label{tab:servicios}
  \begin{tabular}{@{}lp{0.62\textwidth}@{}}
    \toprule
    \textbf{Línea} & \textbf{Alcance} \\
    \midrule
    Desarrollo web & Sitios corporativos y aplicaciones de cliente
      desarrollados a medida \\
    Desarrollo móvil & Aplicaciones multiplataforma \\
    Automatización con inteligencia artificial & Agentes y flujos automatizados
      integrados en los procesos del cliente \\
    Integración \emph{hardware}-\emph{software} & Desarrollo de soluciones que
      articulan componentes físicos con sistemas informáticos \\
    Diseño de interfaz y experiencia de usuario & Definición de la interacción y
      del aspecto visual de los productos \\
    Asesoría tecnológica & Acompañamiento en decisiones de arquitectura y
      selección de herramientas \\
    \bottomrule
  \end{tabular}
\end{table}

De estas líneas, dos resultan directamente pertinentes para el presente trabajo y
se retoman al final del capítulo: la automatización mediante inteligencia
artificial y la integración entre \emph{hardware} y \emph{software}.

\section{Modelo de trabajo}

La empresa acumula más de veinte proyectos completados, con un plazo de entrega
característico en torno a las seis semanas. Esa cifra describe un modo de operar
orientado a ciclos cortos, en los que el producto se acota para poder entregarse
en un plazo acotado en lugar de abordarse como un desarrollo prolongado.

Los proyectos de su portafolio abarcan ámbitos heterogéneos —diseño y arte,
hostelería, interiorismo y comercio electrónico—, lo que indica que la
especialización de la empresa no reside en un sector vertical concreto sino en un
conjunto de capacidades técnicas aplicables a dominios distintos.

La contratación se articula bajo dos modalidades: por proyecto, con un alcance
cerrado de antemano, y por suscripción, que da al cliente acceso continuado a la
capacidad de desarrollo de la empresa. La segunda modalidad supone una relación
sostenida en el tiempo, lo que favorece la acumulación de conocimiento sobre el
dominio de cada cliente.

\section{Ámbito de la pasantía}

El proyecto desarrollado durante esta pasantía se sitúa en la intersección de dos
de las líneas de servicio descritas. Por un lado, se trata de un sistema de
automatización basado en un modelo de lenguaje, esto es, de la aplicación de
inteligencia artificial a una tarea que hasta ahora exigía trabajo manual
especializado. Por otro, su objeto es la generación de firmware para un
microcontrolador, de modo que pertenece al terreno de la integración entre
\emph{hardware} y \emph{software}.

Esa doble pertenencia explica el encaje del proyecto en la empresa. El desarrollo
de firmware exige un conocimiento específico de cada dispositivo que no se
transfiere entre familias de microcontroladores, lo que resulta especialmente
costoso para una organización que trabaja por ciclos cortos y sobre dominios
heterogéneos. Un sistema capaz de producir código verificable a partir de la
documentación del fabricante reduce ese coste de entrada y, en consecuencia, hace
viable abordar proyectos con componente empotrado sin que el plazo de aprendizaje
condicione la planificación.

La supervisión de la pasantía corresponde a Ángel David Rodríguez Niño, Técnico
Superior Universitario en Electrónica y Gerente de Tecnología de la compañía,
designado tutor industrial. La estancia se desarrolla en modalidad presencial en
la sede de la empresa, con una duración de veinte semanas.
